\section{本章小结}
\label{sec:ccode-summary}

本章从容量的存在性结论转向具体的码：在$p=0.1$的二元对称信道上，用有限的块长与可以执行的判决规则纠正翻转，并把信源压缩接到信道传输之前。开头的问题有了答案：$[7,4,3]$ Hamming码只用$7$次信道发送$4$个比特，任意一个翻转都能由综合征直接定位并改正。

\paragraph{主要结论}
\begin{itemize}
 \item 冗余与距离：长度为$2t+1$的重复码在多数判决下恰好纠正至多$t$个翻转，块错误概率由二项分布的尾部给出。一般码本的纠错能力由最小距离$d_{\min}$刻画：保证检测$d_{\min}-1$个翻转，在最小距离译码下保证纠正$\lfloor(d_{\min}-1)/2\rfloor$个翻转，保证来自互不相交的判决球；超出半径后的错误要按具体信道与并列规则计算。

 \item 线性结构：$[n,k]$二元线性码是$\mathbb{F}_2^{n}$的$k$维子空间，最小距离等于最小非零重量。生成矩阵$\bm{G}$把编码变成$\bm{c}=\bm{u}\bm{G}$，存储从$2^{k}n$比特降到$kn$比特；校验矩阵$\bm{H}$满足$\bm{G}\bm{H}^{\top}=\bm{0}$，综合征$\bm{H}\bm{y}^{\top}=\bm{H}\bm{e}^{\top}$只依赖错误。综合征译码选取陪集中重量最小的错误，在二元对称信道上就是最大似然解码；它判对当且仅当真实错误恰为陪集首，但对照表有$2^{n-k}$项，线性结构本身不保证高效译码。

 \item Hamming码：校验矩阵取全部非零列时，各列互异保证纠单错，按二进制排列使综合征直接读出错误位置。它是完备码，在$p=0.1$时块错误$0.150$，并且没有任何$16$条码字、块长$7$的码能做得更好；但对固定的$p$，块长增长时它的错误趋于$1$，不是逼近容量的码族。

 \item 有限块性能：球填充界$MV(n,t)\leq2^{n}$限制保证纠错的码率；在二元对称信道上，$1-P_{\mathrm{e}}\leq g(2^{n}/M)$对任何码与解码器成立；随机编码并集界保证存在达到相应错误的码。$\varepsilon=10^{-3}$、$n=1000$时，最优码率在$0.444$与$0.447$之间，比容量$0.531$低约$16\%$，差距大致按$n^{-1/2}$缩小。

 \item 信源信道分离：$k$个源符号用$n$次信道、$k/n\to r>0$时，无损传输在$rH(S)<C$时可以由先压缩、再用最大错误趋零的信道码实现，任何方案都必须满足$rH(S)\leq C$；有限字母表、有界可分失真下，有损传输的对应条件是$rR(D)<C$与$rR(D)\leq C$。逆定理对任何联合方案成立，所以分离在渐近意义下不损失速率。
\end{itemize}

\paragraph{条件与边界}
距离保证是组合的，与信道概率无关；错误概率则依赖信道、码本和并列规则。线性码与综合征把编码和检错变得便宜，却没有让一般线性码的最大似然译码变得容易。有限块的界限依赖于指定的信道与错误口径，三类陈述（不可能更好、存在这样好、这个码实际多好）不能互相替代。

分离定理只在严格不等式下给出渐近可达性，不承诺有限块、零错误或精确失真；无界失真、有限延迟和带记忆的信源或信道都要另外控制尾部损失、核算有限块代价或补充过程条件。均匀二元信源与高斯信源的不编码传输说明，匹配的联合方案可以在零延迟下达到分离方案只能渐近接近的界限。

\begin{BookInlineTable}
 \centering
 \caption{本章的几个码在$p=0.1$的二元对称信道上的表现}
 \label{tab:ccode-summary}
 \begin{tabular}{lllll}
 \toprule
 码 & 码率 & $d_{\min}$ & 块错误概率 & 比特错误率 \\
 \midrule
 不编码 & $1$ & $1$ & $0.1$（每比特） & $0.1$ \\
 三次重复 & $1/3$ & $3$ & $0.028$ & $0.028$ \\
 五次重复 & $1/5$ & $5$ & $0.00856$ & $0.00856$ \\
 $[5,2,3]$线性码 & $2/5$ & $3$ & $0.067$ & $0.042$ \\
 $[7,4,3]$ Hamming & $4/7$ & $3$ & $0.150$ & $0.067$ \\
 \bottomrule
 \end{tabular}
\end{BookInlineTable}

本部分到此研究了在给定信道上怎样选择输入与添加冗余来可靠传输。下一部分把概率分布本身作为待求的对象：第~\ref{chap:distribution-selection}章研究目标与约束怎样确定所选的分布，其后各章研究怎样逼近一个给定的分布。
